\MathBookSourcePage{387}
\subsection{同时逼近定理}
\label{sec:ch14-simultaneous-approximation-theorem}
\setcounter{thmcount}{0}
\setcounter{equation}{0}

本书中已经多次遇到这种情形, 文献中还有更多例子: 必须知道某个特定逼近元(例如插值函数)同时在两个不同范数中具有适当的逼近阶. 假设~\eqref{eq:ch08-weighted-interpolation-estimates} 就是一个例子. 事实证明, 在许多情形下这种假设是多余的, 因为一个范数中的可逼近性可以推出同时可逼近性. 下面给出一个简单例子, 一般理论见 Bramble 与 Scott (1978).

假设有一族空间 $V_h\subset H^m(\Omega)$, 满足对所有 $u\in H^k(\Omega)$ 和 $0<h\leq1$,
\MBSourceItem{1}
\begin{equation}
\MathBookFormulaID{formula-ch14-single-norm-approximation-assumption}
\label{eq:ch14-single-norm-approximation-assumption}
\inf_{v\in V_h}\|u-v\|_{H^m(\Omega)}
\leq Ch^{k-m}\|u\|_{H^k(\Omega)}.
\end{equation}
下述结果见 Bramble 与 Scott (1978).

\MBSourceItem{2}
\begin{Theorem}
\label{thm:ch14-simultaneous-approximation}
假设条件~\eqref{eq:ch14-single-norm-approximation-assumption} 成立, 且 $s<m$ 和 $m\leq r\leq k$. 则存在常数 $C$, 使得当 $u\in H^r(\Omega)$ 时,
\[
\MathBookFormulaID{formula-ch14-simultaneous-approximation-bound}
\inf_{v\in V_h}\left(h^s\|u-v\|_{H^s(\Omega)}
+h^m\|u-v\|_{H^m(\Omega)}\right)
\leq Ch^r\|u\|_{H^r(\Omega)}.
\]
若 $r<k$ 且 $u\in H^r(\Omega)$, 则
\[
\MathBookFormulaID{formula-ch14-simultaneous-approximation-little-o}
h^{-r}\inf_{v\in V_h}\left(h^s\|u-v\|_{H^s(\Omega)}
+h^m\|u-v\|_{H^m(\Omega)}\right)\longrightarrow0
\quad\text{当 }h\to0.
\]
\end{Theorem}

注意定理中的 $s$ 可以为负, 而且结果容易推广到任意有限个范数之和, 不限于两个. 其他 Sobolev 空间($p\neq2$)的相应结果可由 Bramble 与 Scott (1978) 得到. 不过, 类似~\eqref{eq:ch08-weighted-interpolation-estimates} 的条件并不容易由此推出, 因为它涉及一族加权范数, 且权依赖于逼近参数 $h$.
